% §4 Analysis: Ceasing and Hiding Latency
We state and prove the two-floor results.
Throughout, $\Lambda_{on}$ and $\Lambda_c$ denote the onset and cessation floors of
Definition~\ref{def:floors}.
Constituent latencies of the pipeline are written $\lambda_{\cdot}$, with superscript $irr$
meaning \emph{irreducible}; reaction components of an interruption are written
$\delta_{\cdot}$ (\emph{detect}, \emph{decide}, \emph{cease}, \emph{resume}).
$\tau_f$ is the duration of one synthesis frame and $d$ the committed buffer depth of the
delivery clock.
By a ``floor'' we mean the infimum latency over admissible scheduler decisions under the fixed
in-flight policy of Definition~\ref{def:floors}.

\subsection{The dependency structure of the two floors}
We state the dependency relation as a proposition rather than a theorem: it is close to
definitional, but the shared subsystems are the load-bearing observation.

\begin{proposition}[Floor-constituent dependency]
\label{prop:floor_dependency}
Under a fixed in-flight policy (fixed detection configuration, buffer depth $d$, decode
speed), the onset floor and the cessation floor are set by largely disjoint pipeline
segments, except for three shared subsystems.
Writing $\lambda_s^{on}$ and $\lambda_s^{c}$ for the onset-side and cessation-side
contributions of a shared subsystem $s \in \{ \mathrm{detect}, \mathrm{buf},
\mathrm{decode} \}$, the shared components satisfy
\begin{equation}
  \lambda_{\mathrm{detect}}^{c} = \lambda_{\mathrm{detect}}^{on}, \qquad
  \lambda_{\mathrm{buf}}^{c} = d \cdot \tau_f, \qquad
  \lambda_{\mathrm{decode}}^{c} \leq \lambda_{\mathrm{decode}}^{on},
\end{equation}
where the buffer depth $d$ is the same parameter appearing in the cessation bound of
Theorem~\ref{thm:cessation_floor}.
Consequently, a stage-level optimisation acting only on a single-stage constituent (for
example prefill) moves a single floor; an optimisation acting on a shared subsystem moves
both.
\end{proposition}

\begin{proof}
Onset is gated by detection, prefill/decode, and the buffer handing the first unit to
synthesis.
Cessation is gated by what the delivery and tool clocks still owe after an interruption.
Detection is the same VAD feeding both timelines, giving the equality.
The buffer depth appears in both as the same committed-tail parameter $d \cdot \tau_f$.
Decode contributes to cessation only by filling the committed buffer, which is capped at $d$
frames, hence its cessation-side effect never exceeds its onset-side effect.
\end{proof}

\subsection{The cessation floor}
We decompose interruption reaction into detection, decision, cession, and resumption, and
condition on the interruption phase (Definition~\ref{def:phases}).

\begin{theorem}[Phase-conditioned cessation floor]
\label{thm:cessation_floor}
Let an interruption be detected at time $t_0$, with detection latency
$\delta_{\mathrm{detect}}$ and cancellation-decision latency $\delta_{\mathrm{decide}} \geq
0$.
Under flush-on-cancel delivery, the cession tail in phase-play satisfies
\begin{equation}
  \delta_{\mathrm{cease}}^{\mathrm{phase\text{-}play}} \;\in\;
  [\tau_f,\; d \cdot \tau_f] + \delta_{\mathrm{detect}} + \delta_{\mathrm{decide}},
\end{equation}
so the worst case is bounded below by $\delta_{\mathrm{detect}} +
\delta_{\mathrm{decide}} + d \cdot \tau_f$.
In phase-gen with nothing audible, $\delta_{\mathrm{cease}} \approx 0$ prior to the decision
terms, and the reaction is task-cancel (onset-adjacent).
In phase-tool, nothing is audible, so cessation is vacuous and the meaningful quantity is
$\delta_{\mathrm{resume}}$.
The \emph{can stop} latency (physical) is set by the in-flight granularity $d \cdot \tau_f$;
the \emph{should stop} decision (policy) is separate and may legitimately delay cession for
engagement reasons.
\end{theorem}

\begin{proof}
Detection and decision necessarily precede any cession action, so both add to the reaction.
In phase-play, each committed frame plays to completion under flush-on-cancel, producing a
tail between one frame (only the currently playing frame is committed) and $d$ frames (a full
committed buffer), i.e.\ $\delta_{\mathrm{cease}} \in [\tau_f, d \cdot \tau_f]$.
In phase-gen, no frame is committed, so the tail is empty.
In phase-tool, audible output is absent by definition, so the observable is whether execution
persists, which sets $\delta_{\mathrm{resume}}$.
\end{proof}

\begin{remark}
\label{rem:saturation}
Under the fixed in-flight policy the phase-play tail \emph{saturates} at $d \cdot \tau_f$: a
longer agent utterance lengthens the cease tail only up to the point where the committed
buffer is full.
The general decomposition $\mathrm{reaction} = \delta_{\mathrm{detect}} +
\delta_{\mathrm{decide}} + \delta_{\mathrm{cease}} + \delta_{\mathrm{resume}}$, specialised
to phase-play, composes the cessation floor $\Lambda_c \geq \delta_{\mathrm{detect}} +
\delta_{\mathrm{decide}} + \tau_f$: the minimum stop tail under the fixed policy.
\end{remark}

\begin{remark}
The companion paper's \emph{learned predictive stop} ceiling and our reactive detect--cease
floor together delimit the intervention window: an agent that predicts an upcoming
interruption can begin ceding before detection, contracting $\delta_{\mathrm{detect}}$ toward
its floor.
For long sessions, $\delta_{\mathrm{resume}}$ is additionally bounded below by re-prefill of
the retained context, which the forgetting-floor problem of
Section~\ref{sec:open_problems} targets.
\end{remark}

\subsection{Tool-hiding coupling}
Tools couple the two floors through the onset side unless scheduled onto the tool clock.

\begin{theorem}[Tool-hiding coupling, serial chains]
\label{thm:tool_hiding}
Let a turn require a set of tool calls typed per Definition~\ref{def:tool_typing}, executed
in a serial chain.
Under serial (in-path) dispatch, the onset floor includes the sum of irreducible type-A round
trips $\sum_{i \in A} \lambda_{tool,i}^{irr}$.
Under background dispatch with model reentry on completion
(Definition~\ref{def:interruptibility}), each round trip is removed from the first-audio gate
of the segment it precedes, and only a tool whose result is genuinely required before any
audio of that segment can be produced keeps its type-A cost on the floor.
For a parallel tool DAG, the onset term is the maximum (critical-path) sum over the branches,
not the total over all tools.
\end{theorem}

\begin{proof}
Serial dispatch executes each call in the same path as generation, so first audio of the
chain cannot precede the last type-A result on which the opening utterance depends; the sum
is the chain total.
Background dispatch moves execution to the tool clock, so each completion triggers a fresh
segment whose own onset is gated only by that completion.
For parallel branches the floor is the critical-path maximum, per the companion
paper's DAG latency result.
\end{proof}

The typing here is a \emph{when-needed} classification (does the audio wait on this
result?); it is not a persistence-safety classification, which
Section~\ref{sec:open_problems} takes up.

\subsection{Serializability of tool persistence}
The safety of generation-interruptible, execution-persistent behaviour rests on a commit
invariant.

\begin{lemma}[API serializability as an invariant]
\label{lem:serializability}
For a provider whose function-calling protocol rejects assistant messages with incomplete
tool-call batches, a conversation history is API-valid (Definition~\ref{def:apivalid}) if and
only if every assistant tool-call batch is committed atomically with its complete result set.
Keeping a pending batch in internal orchestrator state (never flushed to the provider) is
safe by construction; the invalid operation is \emph{flushing} a partial batch to the
provider.
\end{lemma}

\begin{proof}
Flushing a partial batch produces an assistant message whose tool-call set lacks its complete
result set, which such protocols reject.
Atomic commit is the only writer path that avoids it; internal retention keeps partial
batches entirely out of the provider's view.
\end{proof}

\subsection{Filler placement}
Filler interacts with both floors through the clock that emits it.

\begin{corollary}[Filler hygiene]
\label{cor:filler}
Masking latency with filler is an \emph{onset} operation and must be emitted by the delivery
clock, never committed to the generation stream.
Short manual filler additionally shortens the phase-play ceasing tail by ending earlier
(Theorem~\ref{thm:cessation_floor}).
Model-uttered filler that is committed to the conversation stream must be rejected on
content-integrity grounds: append-only provider histories cannot retract a partial filler
fragment, and unguarded filler bleeds into or duplicates the committed response.
\end{corollary}

\begin{proof}
Filler that precedes the first real utterance shortens perceived onset by
Definition~\ref{def:floors}; a manually played clip is delivered and forgotten, never
entering history.
On append-only histories a committed model-uttered fragment is unretractable, so an
interrupted filler corrupts the committed record; reliability of the delivered text therefore
forces filler to stay off the generation stream.
\end{proof}

\subsection{Smoke anchor}
We instantiate the semantics with a small event-driven orchestrator
(generation-interruptible, execution-persistent, atomically committed tool batches) and
anchor the quantitative claims with a monotonic-clock delivery-clock simulation.
Table~\ref{tab:smoke} reports five-run medians.

\begin{table}[t]
\centering
\begin{tabular}{lr}
\toprule
\textbf{Metric} & \textbf{Median (ms)} \\
\midrule
cease, buffer depth 1 & 109 \\
cease, buffer depth 6 & 609 \\
generation-stop, depth 1 / 6 & $<$0.1 / $<$0.1 \\
resume, tool 10 ms (in-memory) & $<$1 \\
resume, tool 1.0 s & 1000 \\
\bottomrule
\end{tabular}
\caption{Smoke anchor of the cessation floor (five-run medians, monotonic clock, simulated
delivery clock with committed buffer depth $d$). The scripted model has zero prefill cost, so
resume figures bound only reentry plus remaining tool delay. An anchor, not a benchmark:
there is no real ASR/TTS/LLM, and frame and tool latencies are synthetic.
\label{tab:smoke}}
\end{table}

In phase-play the measured cease is linear in the committed buffer depth
($109$~ms at $d=1$ vs $609$~ms at $d=6$ for a $100$~ms synthetic frame), consistent with
Theorem~\ref{thm:cessation_floor}'s $[\tau_f, d \cdot \tau_f]$ bound; in-fill is synthetic,
so the numbers are a self-consistent reproduction, not independent field evidence.
Generation cancellation is prompt ($<$0.1~ms median) and depth-independent, as the decoupled
``cancel generation, persist tools'' semantics predict.
In phase-tool the model task has already completed, so the interrupt call cancels nothing new;
the measured resume therefore records the tool-completion reentry path, and it is dominated by
the tool round trip ($< 1$~ms in-memory vs $\sim$1000~ms for the $1.0$~s tool), consistent
with Theorem~\ref{thm:tool_hiding}.

\begin{figure}[t]
\centering
% The two responsiveness floors on a single turn timeline.
\begin{tikzpicture}[>=Stealth, font=\small]
  % Timeline axis
  \draw[thick,->] (0,0) -- (12,0);
  % Marks
  \draw (1,0) -- (1,-0.12); \node[below,align=center] at (1,-0.5) {end of \\ user speech};
  \draw (6,0) -- (6,-0.12); \node[below,align=center] at (6,-0.5) {first agent \\ audio};
  \draw (8.5,0) -- (8.5,-0.12); \node[below,align=center] at (8.5,-0.5) {interrupt \\ detected};
  \draw (11.5,0) -- (11.5,-0.12); \node[below,align=center] at (11.5,-0.5) {fully \\ stopped};
  % Agent audio bar
  \draw[very thick,blue] (6,0) -- (8.5,0);
  \node[blue,above] at (7.25,0.15) {agent audio};
  % Onset span
  \draw[<->,green!50!black] (1,0.7) -- (6,0.7);
  \node[green!50!black,above] at (3.5,0.9) {onset floor $\Lambda_{on}$};
  % Cessation span
  \draw[<->,red] (8.5,0.7) -- (11.5,0.7);
  \node[red,above] at (10,0.9) {cessation floor $\Lambda_c$};
  % Shared subsystems note (box sized to the two text lines)
  \draw[dashed] (0.4,-1.7) rectangle (11.6,-2.6);
  \node[align=center] at (6.0,-2.15) {shared by both floors:\\ detection $\cdot$ buffer policy $\cdot$ decode speed};
\end{tikzpicture}
%
\caption{The two responsiveness floors on a single turn timeline. Onset spans the detectable
end of user speech to first audible agent audio; cessation spans interruption onset to the end
of all in-flight audio. The detection, buffer, and decode subsystems are shared by both floors.
\label{fig:two_floors}}
\end{figure}

\begin{figure}[t]
\centering
% The three interruption phases and the react path each one takes.
\begin{tikzpicture}[>=Stealth, font=\small, node distance=9mm and 14mm,
  box/.style={draw, rounded corners, align=center, minimum height=8mm}]
  % Phase 1
  \node[box] (pgen) {phase-gen\\nothing audible};
  \node[box, right=of pgen] (rg) {cancel generation task};
  \node[right=of rg] (og) {$\delta_{cease} \approx 0$};
  \draw[->] (pgen) -- (rg); \draw[->,gray] (rg) -- (og);
  % Phase 2
  \node[box, below=of pgen] (ppl) {phase-play\\audio delivered};
  \node[box, right=of ppl] (rpl) {flush delivery clock\\ in-flight tail plays};
  \node[right=of rpl] (opl) {$\delta_{cease} \approx d\cdot\tau_f$};
  \draw[->] (ppl) -- (rpl); \draw[->,gray] (rpl) -- (opl);
  % Phase 3
  \node[box, below=of ppl] (ptl) {phase-tool\\tool waiting};
  \node[box, right=of ptl] (rtl) {execution persists\\resume on completion};
  \node[right=of rtl] (otl) {$\delta_{resume}$};
  \draw[->] (ptl) -- (rtl); \draw[->,gray] (rtl) -- (otl);
\end{tikzpicture}
%
\caption{The three interruption phases and the react paths each one takes: phase-gen is
task-cancel, phase-play is delivery flush, phase-tool reduces to resumption after the
background tool completes.
\label{fig:interruption_phases}}
\end{figure}

\begin{figure}[t]
\centering
% Event algebra of the orchestrator: generic events, three clocks, reentry.
\begin{tikzpicture}[>=Stealth, font=\small,
  st/.style={draw, rounded corners, align=center, minimum height=9mm, minimum width=24mm},
  ev/.style={->,>=Stealth}]
  \node[st] (idle) at (0,0) {Idle};
  \node[st] (gen)  at (5,1.5) {Generating};
  \node[st] (del)  at (10,1.5) {Delivering};
  \node[st] (tool) at (5,-2.5) {Tool running};
  % user turn: idle -> generating (diagonal)
  \draw[ev] (idle) -- (gen) node[midway,above,font=\scriptsize,left=2mm]{user turn};
  % stream: generating -> delivering (straight, label above)
  \draw[ev] (gen) -- (del) node[midway,above,font=\scriptsize]{stream};
  % interrupt: red arc OVER the top left (cancels generation only, not tools)
  \draw[ev,red] (gen.north west) to[bend left=75]
    node[midway,above,font=\scriptsize,red]{interrupt} (idle.north east);
  % tool call: right bulge down from generation to the tool
  \draw[ev] (gen.south east) .. controls (6.3,-0.3) and (6.3,-1.5) .. (tool.north east)
    node[pos=0.5,right,font=\scriptsize]{tool call};
  % completion -> reentry: left bulge up from the tool to generation
  \draw[ev] (tool.north west) .. controls (3.7,-1.6) and (3.7,-0.2) .. (gen.south west)
    node[pos=0.5,left,font=\scriptsize]{completion $\rightarrow$ reentry};
  % tool execution persists across interruption (self-loop under the tool)
  \draw[dashed,->,gray] (tool.south) .. controls (6.4,-3.3) and (3.6,-3.3) .. (tool.south);
  \node[gray,font=\scriptsize] at (5.15,-3.55) {execution persists};
  % delivery outcomes: one outer loop around the right and bottom (clear of everything)
  \draw[ev] (del.south) .. controls (11.7,-3.5) and (0.3,-4.6) .. (idle.south)
    node[pos=0.30,below,font=\scriptsize,xshift=1.5mm,yshift=-1mm]{delivery: completed / interrupted};
\end{tikzpicture}
%
\caption{The event algebra of the orchestrator: user turns and tool completions are generic
events; delivery status reports playback; interruption cancels generation while the tool clock
persists and re-enters on completion.
\label{fig:event_loop}}
\end{figure}

\begin{figure}[t]
\centering
% Tool placement against the onset floor: only a serial type-A tool gates first audio;
% background dispatch, lead-in hiding, and background clocks all move off the floor.
\begin{tikzpicture}[>=Stealth, font=\small, box/.style={draw, rounded corners,
  align=center, minimum height=7mm}]
  % Onset floor band (top, horizontal)
  \draw[fill=green!8, draw=green!35!black, dashed, rounded corners] (-2.2,3.2) rectangle (9.6,4.2);
  \node[green!35!black] at (3.5,3.7) {onset floor: first-audio gate};
  % Serial type A -> INTO the band (gates)
  \node[box] (as) at (0,1.9) {type A, serial dispatch};
  \draw[->,thick,green!35!black] (as.north) -- (1.4,3.2) node[midway,right,font=\scriptsize]{gates};
  % Type A background -> away (down)
  \node[box] (ab) at (0,0.7) {type A, background dispatch};
  \draw[->,thick] (ab.south) -- (0,-0.5) node[midway,right,font=\scriptsize]{removed from floor};
  % Type B -> away (down)
  \node[box] (bt) at (4.2,0.7) {type B, deferrable};
  \draw[->,thick,gray] (bt.south) -- (4.2,-0.5) node[midway,right,font=\scriptsize]{hidden under lead-in};
  % Type C -> away (down)
  \node[box] (ct) at (7.6,0.7) {type C, later turn};
  \draw[->,thick,gray] (ct.south) -- (7.6,-0.5) node[midway,right,font=\scriptsize]{background tool clock};
\end{tikzpicture}
%
\caption{Tool placement against the onset floor: a serial type-A tool gates first audio
(inside the band); background dispatch removes it from the floor (arrow away from the band);
type-B is hidden behind delivery lead-in and type-C behind the background tool clock.
Only the earliest genuinely audio-gating type-A result stays on the floor; for a parallel
DAG it is the critical path, not the sum.
\label{fig:tool_hiding}}
\end{figure}
